% !PRINTED 522
\weberpage{522}

Here there now stand, on the right hand and on the left, integral functions of $t$ whose coefficients are integral functionals. On the right side all these coefficients have the factor $\mu$, hence also the factor $\pi$, while by assumption not all the coefficients of $\chi$ have this factor. Since $E$ is a unit, the coefficients of $E$ likewise do not contain the factor $\pi$, and consequently, by Theorem 1., the coefficients of $E\chi^m$ also cannot all be divisible by $\pi$, which the equation (4) would nevertheless require. From this it follows that our assumption, that $\varphi$ is integral while $\varphi_0, \varphi_1, \ldots, \varphi_h$ by contrast are not all integral, is inadmissible, and Theorem 2. is thus proved.

If we now assume that in $\varphi$ the coefficients $\varphi_0, \varphi_1, \ldots$ are themselves in turn integral functions of a variable, and apply Theorem 2. repeatedly to this, we arrive at the conclusion:

3. \gesperrt{An integral rational function of arbitrarily many variables, whose coefficients are numbers or functionals with other variables, is an integral functional only when the coefficients are integral.}

We can represent every functional $\omega$ as a quotient of two integral functions in such a way that the denominator becomes a primitive function in the field $R$.

For if $\varphi, \psi$ are integral functions in $\Omega$, and if

\[
\omega = \frac{\varphi}{\psi}, \quad N(\psi) = \psi\psi',
\]

then we can multiply both terms of the fraction $\omega$ by $\psi'$ and obtain in the denominator an integral function with rational coefficients. We can then reckon the divisor of this function to the numerator of $\omega$, and we obtain, if $E$ denotes a primitive function and $\chi$ an integral function in $\Omega$:

\[
E\omega = \chi,
\]

that is, one can transform every functional $\omega$ in $\Omega$, by multiplication with a primitive function $E$, into an integral function of the variables whose coefficients are numbers in $\Omega$.

With the aid of Theorem 3. we obtain from this:

4. \gesperrt{Every integral functional $\omega$ in the field $\Omega$ is associated with an integral function $\chi$, whose coeffi-}
%CONT>>
